\documentclass[11pt]{article}
\usepackage[utf8]{inputenc}
\usepackage[T1]{fontenc}
\usepackage{lmodern}
\usepackage[margin=1in]{geometry}
\usepackage{csquotes}
\usepackage[style=authoryear,backend=biber,sorting=nyt,maxbibnames=99,giveninits=true,uniquename=init]{biblatex}
\usepackage[hidelinks]{hyperref}
\addbibresource{references.bib}
\begin{document}
\section*{Salinity, Reclamation, and Drainage-Water Reuse in Arid and Irrigated Systems}

Subsurface drainage is a cornerstone technology for reclaiming salt-affected
land, yet its salt-leaching outcomes accrue slowly, are only partially
observable with current programmes, and generate salt-laden water whose disposal
remains unsolved. The frontier here is defined by this triad---slow and
incompletely measured reclamation, the management of the water that reclamation
produces, and the specialised chemistry of acid sulfate and irrigation-return
systems---together with the sensing, modelling, and design tools needed to
resolve it.

\subsection{Salt-leaching dynamics and long-term reclamation trajectories}

The recurring message is temporal: meaningful salt reduction often accrues so
slowly that short observation windows systematically underestimate long-term
benefit, and only extended drainage with sustained monitoring resolves the true
trajectory \parencite{Almen2021,Daigh2025}. The long-term effects of combined
saline-water irrigation and subsurface drainage on both salinity and yield
require further evaluation \parencite{Feng2017,Feng2021}, desalination rates and
trends remain incompletely quantified \parencite{Guo2024b,Yao2024}, and
controlled drainage introduces unresolved temporal trade-offs, buffering
outflow while risking salt accumulation
\parencite{Jouni2018,Lu2019,Islam2025,Heng2021}. Longer observation records
paired with optimal leaching quotas and cost--benefit analysis are needed to
close these gaps \parencite{Guo2025}, and salt-tolerant crops in reclamation
rotations, alongside quantification of rainfall and irrigation-efficiency gains
relative to salt-removal capacity, remain to be characterised
\parencite{Qadir2001,Castanheira2010}.

A parallel gap is measuring what actually leaves through the drains. Priorities
include greater attention to surface-water monitoring of salt and fertiliser
export \parencite{Mastrocicco2013}, continuous flow meters and
electrical-conductivity probes for oasis-scale balances with evaluation of
impacts on downstream phreatic groundwater \parencite{Askri2022}, nutrient
balances of recipient waters \parencite{Turtola1995}, and cluster wells with
deep sampling to track salinity over time \parencite{Singh2020d,Gafni2001}.
Process understanding of salt mobilisation is itself an active frontier: how
macropores govern transient-state leaching and how their effects might be
minimised \parencite{Xin2016}, the controlling factors in low-permeability and
structured media \parencite{Zhang2022}, salt replacement between pipes under
brackish-water leaching \parencite{Li2023}, rainfall--irrigation interactions
and generalisable desalinisation ratios across climates and soils
\parencite{Yu2015,Chen2025,Wang2024d}, ponding strategies and leaching curves
\parencite{RAO1991,Yang2022c}, and the optimal root-zone matric potential that
links leaching intensity to crop response \parencite{Zhang2021,Heng2021}.
Overarching these, waterlogging and salinity should be assessed together rather
than separately \parencite{Valipour2014}, and the economic feasibility of
drainage as a preventive measure against salinisation remains to be established
\parencite{Datta2000,HalbacCotoaraZamfir2022}.

\subsection{Drainage-water reuse, disposal, and treatment}

Removing salt from soil relocates the problem to the drainage water, whose
management is among the most persistent open challenges. Environmentally
acceptable disposal methods are still needed \parencite{Ayars2014,Steppuhn2006},
evaporation-pond management must be tested at full scale and combined with
volume- and contamination-reduction measures \parencite{BRADFORD1991}, and
trace-element fate---particularly arsenic speciation and uranium recovery or
removal---compounds the disposal burden \parencite{Gao2007,Duff1997,Gao2003}.
Emerging treatment routes remain at the demonstration stage and require testing
under realistic conditions: electrolysis of saline wastewater at varying
salinity \parencite{Yang2023}, sorptive filter materials evaluated with real
rather than artificial drainage water \parencite{Pugliese2023}, and blending of
drainage water with reverse-osmosis product water \parencite{Singh2020d}.
Life-cycle cost assessment, scalable data infrastructure, and low-temperature
crystallisation toward near-zero discharge define the longer-horizon disposal
agenda \parencite{Qin2025}.

Reuse research spans site-specific load--flow relationships, germination-stage
salinity tolerance, and improved scheduling that meets crop evapotranspiration
from a shallow water table while minimising upward salt flow
\parencite{WESTCOT1988}. Additional needs include provisional field guidance for
substituting drainage water on wheat \parencite{Sharma1998}, blending combined
with evaporation-pond storage \parencite{Datta2002}, farm-level reuse embedded in
a basin-wide salt balance \parencite{Qureshi2008}, exploitation of seasonal
groundwater undulation and drip conversion to prevent accumulation
\parencite{Mirlas2021,Dou2022,Singh2019}, sequential reuse across cereals,
grasses, forest plantations, and halophytes with innovative evaporation-pond
designs \parencite{Sharma2004}, and land-treatment/pond combinations under
restricted hydraulic loading with preferential-flow control and pre-treatment
before wastewater application \parencite{Christen2010,Duan2015}. Biochar,
phytoremediation, and related reclamation innovations warrant further
development \parencite{Liu2025f,Steppuhn2017,Youngs2000}. Crucially, economic
analysis indicates that feasible drainage-water recycling requires markedly
higher and less volatile irrigated yields together with lower adoption costs
than currently prevail \parencite{Sellars2022}.

\subsection{Acid sulfate soils and irrigation-return management}

Drained acid sulfate (AS) and boreal soils raise a distinct set of specialised
questions. Short-term hydro-meteorological controls on sulfide oxidation,
ripening effects on hydraulic properties, and secondary-mineral formation should
be incorporated into model applications \parencite{Salo2023}; the role of the
nitrogen cycle in sulfide oxidation is unknown in boreal cultivated AS soils
\parencite{Virtanen2017}; the peat-cover thickness required to hinder oxidation
of sulfidic material needs detailed study \parencite{YliHalla2022}; and more
detailed AS-soil maps, a clearer distinction between anthropogenic and organic
acidity, and cessation of drainage of potential AS soils are urged
\parencite{Åström2005}. Amendment strategies---gypsum, elemental sulphur, and
lime for long-term management of problem soils \parencite{MANN1982,He2015},
biochar selected for low contaminant content with urgent long-term impact
assessment \parencite{Tian2023b}, and optimum rates of organic conditioners that
minimise adverse water-quality impacts \parencite{Zhang2015}---require
longer-term verification.

For reclaimed saline, waterlogged, and Vertisol systems, priorities include how
subsurface drainage integrated with crop and residue management affects soil
quality, carbon sequestration, nutrient budgets, and greenhouse-gas fluxes
\parencite{Mukhopadhyay2023,Mukhopadhyay2023b}, the long-term effects of
controlled drainage on salinity \parencite{ElGhannam2023}, salt-crust formation
and dissolved-species transport \parencite{Colombani2016}, leaching effects on
additional elements and organic modifiers \parencite{SoltaniGerdefaramarzi2020},
monitoring of elemental toxicities with waterlogging-stage-specific drainage
scheduling \parencite{Qian2021}, long-term root-growth resistance and
drainable-porosity determination \parencite{Bahçeci2018,Bahceci2008},
additive-effectiveness and land-replacement characterisation \parencite{Isa2023},
and multi-location, multi-condition testing over longer periods
\parencite{Sadhukhan2019}. On the irrigation-return side, needs span tolerant
rootstocks, halophyte germplasm, and transgenic crops \parencite{Cuevas2019},
land-use planning that steers irrigation toward low-salinity aquifers
\parencite{Bethune2008}, and sediment- and return-flow reduction from
furrow-irrigated fields with improved pond removal and nitrogen-leaching-limiting
management \parencite{Bjorneberg2020}. Systematic evaluation of alternative
field-scale irrigation and drainage strategies under ambiguity, expressed as a
statistical likelihood of achieving goals, is advanced \parencite{Alzraiee2013},
supported by long-term geohydrological and crop data to validate simulations
across irrigation/drainage systems \parencite{Yao2017}. Retired land over shallow
water tables, surrounded by irrigated fields, poses unresolved ultimate-use
questions \parencite{Purkey2001}.

\subsection{Sensing, modelling, and design under salinity}

Salinity-specific sensing is a fast-growing cluster: EMI mapping validated across
soil orders and settings \parencite{Narjary2024}, coupled
apparent-conductivity--water-content models \parencite{Yu2025}, and EM38 with
multispectral and hyperspectral remote sensing integrated through interfaces that
explicitly incorporate tile drains \parencite{Singh2020d,Jayawardane2001}.
Coupled water--salt transport modelling remains immature. Models are needed that
simulate salt and nitrate simultaneously in drainage water and the soil profile
\parencite{Ebrahimian2015}; neural networks can predict effluent salinity from
minimal input, but process codes such as SALTMOD may be preferred for detailed
salt balances \parencite{Sarangi2006}, and further evaluation must precede firm
salt-transport recommendations \parencite{Hirekhan2007}. Several strands press
for tighter coupling of water and salt dynamics with unsaturated-zone models such
as HYDRUS, SWAP, and UBMOD-S, incorporating regional layout parameters and
multi-objective optimisation under climatic, textural, and groundwater
variability \parencite{Han2025,Wang2025c,Wang2024c}, and for root-growth modules,
automatic TDR monitoring, and large-dataset validation
\parencite{Liu2022}, extending to climate-change scenarios, deep-soil salt
accumulation, and real-world scalability \parencite{Liu2026}. Sustainability of
saline-water irrigation warrants coupled HYDRUS--EPIC evaluation
\parencite{Feng2019} and better representation of water-table contributions to
storage and evapotranspiration \parencite{Jia2017}. Earlier salinity series
staged validation and salt-flow analysis across successive papers
\parencite{KARAJEH1994}, and proper representation of regional seepage would
demand a fundamentally different modelling approach \parencite{Kelleners2000}.

Design under heterogeneity is the final gap. Integrated redesign of system,
surface-water, and nutrient management is required \parencite{vandenEertwegh2006},
together with texture-specific understanding of the economics of preferential
subsurface drainage \parencite{Neha2022}, and vertical porous flow pipes or fully
charged segments to enhance salt drainage
\parencite{Wang2024d,Yao2024,Jayawardane2001}. New design criteria that protect
crops while minimising both the salinity and the volume of drain water are sought
\parencite{Christen2001}. The most integrative vision unifies these threads:
zonal design algorithms driven by a multi-scale water-and-salt transport model
coupled to remote-sensing inversion for real-time early warning
\parencite{Zhu2025b}---a target that ties salinity sensing, coupled modelling, and
adaptive design into a single research programme.

\printbibliography
\end{document}
